% !TeX program = xelatex
% 适配 zju-resume-template 仓库的 XeLaTeX 简历源文件
% 使用方法：将本文件替换仓库根目录下的 CV.tex，然后由 GitHub Actions 自动编译。
% 请将下方方括号中的个人信息替换为真实信息；不要保留未填写的占位符。

\documentclass[11pt,a4paper]{article}

% -------------------- 基础设置 --------------------
\usepackage{fontspec}
\usepackage{xeCJK}
\usepackage{xltxtra}
\usepackage{hyperref}
\hypersetup{hidelinks}
\usepackage{url}
\urlstyle{tt}
\usepackage{xcolor}
\usepackage{calc}
\usepackage{graphicx}
\usepackage{enumitem}
\usepackage{titlesec}
\usepackage{fontawesome}
\usepackage{microtype}

\setlength{\parindent}{0pt}
\setlength{\parskip}{0pt}
\CJKsetecglue{}
\linespread{1.13}

% 模板仓库内置字体。注意中文字体文件实际扩展名为 .ttf。
\setmainfont[
    Path = fonts/Main/,
    BoldFont = texgyretermes-bold.otf
]{texgyretermes-regular.otf}

\setCJKmainfont[
    Path = fonts/hansans/,
    BoldFont = NotoSansSC-Bold.ttf
]{NotoSansSC-Regular.ttf}

% -------------------- 页面与视觉样式 --------------------
\definecolor{CVBlue}{RGB}{23,110,191}
\usepackage[
    a4paper,
    left=1.35cm,
    right=1.35cm,
    top=1.15cm,
    bottom=1.15cm,
    nohead
]{geometry}

\setlist[itemize]{
    leftmargin=1.45em,
    itemsep=0.12em,
    topsep=0.18em,
    parsep=0pt,
    partopsep=0pt
}

\titleformat{\section}
  {\large\bfseries\raggedright}
  {}{0em}{}
  [{\color{CVBlue}\titlerule}]
\titlespacing*{\section}{0pt}{0.85em}{0.45em}

\newlength{\projectSpacing}
\setlength{\projectSpacing}{0.48em}
\newcommand{\projectsep}{\vspace{\projectSpacing}}
\newcommand{\projecttitle}[2]{%
  \textbf{#1}\hfill{\small\color{black!65}#2}\par\vspace{0.12em}%
}
\newcommand{\cvsection}[2]{%
  \section{\texorpdfstring{\makebox[\widthof{#1}][c]{\color{CVBlue}#1}\hspace{0.35em}#2}{#2}}%
}

\begin{document}
\pagenumbering{gobble}

% -------------------- 个人信息 --------------------
\begin{center}
    {\LARGE\bfseries [姓名]}\\[0.28em]
    {\small
      \faPhone\ [联系电话]
      \quad
      \faEnvelopeO\ \href{mailto:your.email@example.com}{your.email@example.com}
      \quad
      \faGithubSquare\ \href{https://github.com/lilinlilin909-ux}{github.com/lilinlilin909-ux}
    }\\[0.28em]
    {\normalsize\bfseries 求职意向：人形机器人运动控制算法实习生}
\end{center}

% -------------------- 教育背景 --------------------
\cvsection{\faGraduationCap}{教育背景}
\textbf{[硕士就读院校]}\hfill [入学年月] -- 至今\\[-0.05em]
硕士研究生\quad 车辆工程\\[-0.05em]
研究方向：无人电动摩托车智能驾驶

% -------------------- 个人简介 --------------------
\cvsection{\faInfo}{个人简介}
车辆工程专业硕士研究生，研究方向为无人电动摩托车智能驾驶。具备复杂机械系统非线性多刚体动力学建模、现代控制理论及电机控制仿真研究经历。能够从物理原理出发开展动力学方程推导、状态空间建模、控制器数学设计与稳定性分析。独立完成 Whipple 自行车四刚体动力学模型推导，并通过 Meijaard、Basu 相关 Benchmark 验证；另有 PMSM 双闭环 FOC/滑模观测器仿真建模及 Linux/C++ 电机 CAN 通信项目实践。希望将已有动力学与控制研究能力迁移至机器人运动控制领域。

% -------------------- 科研与项目经历 --------------------
\cvsection{\faUsers}{科研与项目经历}

\projecttitle{Whipple 自行车非线性多刚体动力学建模与平衡控制}{[起止时间]}
\begin{itemize}
    \item 基于物理原理建立 Whipple 自行车四刚体动力学模型，采用牛顿—欧拉方法推导刚体动力学方程及运动约束关系，完成动力学公式的显式展开。
    \item 参考 Meijaard、Basu 相关经典研究中的建模方法与推导技巧，完成模型对照验证，并通过 Meijaard 与 Basu 相关的全部 Benchmark。
    \item 基于动力学模型建立状态空间表达，独立开展控制器数学设计、稳定性分析与参数设计，实现自行车转向跟踪及原地平衡仿真。
    \item 模型范围采用经典 Whipple 刚体拓扑，未纳入悬架、制动及轮胎变形等扩展因素。
\end{itemize}

\projectsep
\projecttitle{PMSM 双闭环 FOC 与滑模观测器仿真建模}{[起止时间]}
\begin{itemize}
    \item 基于 MATLAB/Simulink 搭建永磁同步电机（PMSM）双闭环磁场定向控制（FOC）仿真模型，建立电流环与速度环控制结构。
    \item 搭建滑模观测器（SMO）仿真模型，开展电机状态观测相关方法的建模与分析。
    \item 了解基于 TI 电机控制平台的双环参数整定方法；相关经历以 MATLAB/Simulink 仿真建模为主。
\end{itemize}

\projectsep
\projecttitle{RS10P 电机私有 CAN 通信软件开发}{[起止时间]}
\begin{itemize}
    \item 基于 Linux/Ubuntu、C++ 与 SocketCAN，开展 RS10P 电机厂商私有 CAN 2.0 扩展帧协议通信程序开发。
    \item 根据厂商协议手册开展通信帧收发、协议数据解析及控制指令交互等功能实现与调试。
    \item 使用 AI 辅助代码实现、中文注释生成、代码阅读与调试；本人负责协议理解、控制逻辑设计及测试结果分析。
\end{itemize}

% -------------------- 专业技能 --------------------
\cvsection{\faCogs}{专业技能}
\begin{itemize}
    \item \textbf{动力学与运动学：} 牛顿—欧拉法、多刚体动力学建模、刚体运动学、运动约束分析；具备李群与李代数相关理论基础。
    \item \textbf{控制理论：} 经典与现代控制理论；LQR、MPC、$H_\infty$、$H_2$ 控制理论基础；控制器设计、稳定性分析与参数整定。
    \item \textbf{系统辨识：} ARX、ARMAX、状态空间（SS）辨识及 PRBS 激励方法。
    \item \textbf{电机控制：} PMSM 数学建模、FOC 双闭环控制、滑模观测器（SMO）及 TI 平台双环参数整定方法。
    \item \textbf{开发与仿真：} MATLAB/Simulink；Linux/Ubuntu；C++、SocketCAN 电机通信项目实践；AI 辅助代码实现与调试。
\end{itemize}

\end{document}
